\section{Introduction}\label{sec:introduction}

Actuarial decisions under stress require more than an estimate of expected loss. Insurers and regulators also need to know how adverse conditions interact over time, which combinations cross operational or solvency thresholds, whether a critical result can be replayed, and which assumptions or events are associated with failure. Conventional frequency-severity models and Monte Carlo simulations remain essential for loss projection, but the surrounding workflow often treats scenario construction, validation, random-number governance, and diagnostic attribution as separate activities.

This separation creates several risks. A simulated outcome can be numerically extreme but structurally inconsistent. A scenario can be labelled severe and plausible without a quantitative plausibility criterion. A result can be reproducible under one execution order yet change after an unrelated component consumes random numbers. A compact point error can coexist with poor predictive-interval coverage. Finally, a failure record can identify the threshold that was crossed without establishing the real-world cause of that failure. These properties (loss projection, scenario plausibility, computational reproducibility, predictive calibration, and attribution) must therefore be evaluated separately.

ARCCRAFT is proposed as an integration architecture for this problem. It represents a synthetic environment as a versioned, time-dependent world; allocates independent pseudorandom streams to world generation, claims, and behaviour; validates initial states through explicit pass, repair, or reject decisions; simulates multi-period outcomes; and records failed worlds in a seed-addressable Failure Atlas. The framework also includes a local reverse-stress search. These components build on established work rather than constituting a new actuarial theory.

The research question is: \textbf{Does integrating explicit multi-period transitions, structural validation, separated random streams, replay semantics, and failure attribution produce measurable improvements in reproducibility, diagnostic information, failure detection, or empirical prediction relative to transparent comparators, and at what cost?} The evaluation is deliberately falsifiable. ARCCRAFT may add governance value without improving predictive accuracy, and individual architectural components may have little effect under the current grammar.

This study makes four bounded contributions. First, it specifies the implemented architecture and maps its claims to executable functions and versioned outputs. Second, it audits a complete motor-insurance table and implements matched aggregate and policy-level actuarial comparators over two temporal origins. Third, it tests local deterministic replay, stream isolation, validator policy, state transitions, and runtime through controlled ablations. Fourth, it reports negative findings, including interval misses, decile miscalibration, lack of policy-severity superiority, and dominance of the traceability gate, rather than defining validation from favourable point estimates.

\section{Literature review and hypotheses}\label{sec:literature-review-and-hypotheses}

\subsection{Actuarial loss modelling and the limits of static projection}\label{sec:actuarial-loss-modelling-and-the-limits-of-static-projection}

Non-life insurance models commonly decompose aggregate loss into claim frequency and claim severity. This decomposition is attractive because claim counts contain many zeros, positive claim amounts are typically right-skewed, and the two components may respond differently to rating and environmental variables. In its simplest form, a collective-risk model generates a claim count and then draws claim amounts, often under conditional independence. Such models remain useful transparent baselines, but the independence assumption is not innocuous when the number and average size of claims share unobserved risk factors or evolve jointly over time.

Several studies have relaxed this assumption. Garrido et al. \cite{Garrido2016} incorporate dependence by including claim count in a conditional severity model and show how the resulting pure premium differs from the independent frequency-severity formulation. Shi et al. \cite{Shi2015} develop hurdle and copula constructions for policy-level frequency and severity, while Frees et al. \cite{Frees2016} extend the problem to multivariate insurance outcomes. Lee and Shi \cite{Lee2019} add longitudinal dependence, allowing past claim frequency and severity to inform the future predictive distribution. Jeong and Valdez \cite{Jeong2020} similarly use dependent random effects to construct predictive compound-risk models, and Gao and Li \cite{Gao2023} link waiting time and average severity to obtain a predictive aggregate-loss distribution. Together, these studies establish that a credible benchmark for a multi-period actuarial framework should not automatically assume that frequency, severity and their temporal dynamics are independent.

The implications for stress testing are twofold. First, a joint frequency-and-severity shock is not equivalent to two unrelated marginal multipliers when the components share latent or observed dependencies. Second, a multi-period stress environment cannot be evaluated adequately by comparing only its terminal expected loss with a historical mean. The temporal ordering of shocks, transitions and behavioural responses may affect both the path to failure and the distribution of the terminal outcome. A procedural stress architecture must therefore demonstrate what its transition rules add beyond a transparent longitudinal frequency-severity model rather than treating temporal sequencing as intrinsically superior.

Flexible actuarial models provide additional comparators. Klein et al. \cite{Klein2014} use Bayesian generalized additive models for location, scale and shape to let covariates affect more than the conditional mean. Modern tree and boosting methods can capture nonlinearities and interactions, but their predictive gains must be balanced against interpretability and calibration. These methods are relevant to ARCCRAFT because a procedural environment could appear useful simply because it embeds nonlinear dependence that a weak baseline omits. Comparator selection must therefore include at least one transparent conventional model and, where the data support it, one dependence-aware or flexible alternative.

\subsection{From predefined shocks to systematic stress scenarios}\label{sec:from-predefined-shocks-to-systematic-stress-scenarios}

Insurance stress testing is intended to examine resilience under adverse conditions, not merely to reproduce the central predictive distribution. EIOPA \cite{EIOPA2020} frames insurance stress testing around clearly specified objectives, material risks and severe-but-plausible scenarios. This supervisory framing makes scenario design a substantive modelling decision: an adverse outcome is informative only if the scenario producing it is coherent with the exercise's purpose and its plausibility criterion. A large loss generated by an internally inconsistent or unconstrained combination of variables is not automatically a credible stress-test result.

The scenario-selection literature formalizes this problem. Breuer et al. \cite{Breuer2009} characterize the tension between severity and plausibility and propose searching systematically for damaging scenarios inside an admissible region. Breuer and Csiszar \cite{Breuer2013} use entropic plausibility constraints to construct systematic stress tests, while Glasserman et al. \cite{Glasserman2015} use empirical likelihood to identify plausible combinations of factors associated with extreme losses. Flood and Korenko \cite{Flood2015} position systematic scenario selection between hand-picked shocks and reverse stress testing, emphasizing internally consistent multidimensional shocks and the reduction of blind spots.

These approaches reveal an important distinction for ARCCRAFT. A structural validator can reject impossible combinations or repair violated bounds, but internal consistency is not identical to statistical or economic plausibility. A world can satisfy coded invariants while remaining implausible relative to observed joint dynamics. Conversely, a repair operation may create a formally admissible world while altering the scenario distribution in ways that bias tail estimates. The validator must therefore be evaluated as part of the data-generating process. Its acceptance, repair and rejection rates, and the resulting changes in marginal and joint output distributions, are empirical outcomes rather than implementation details.

The literature also cautions against presenting a fixed stress label as self-validating. Terms such as \emph{adverse} and \emph{severe plausible} need quantitative definitions: probability regions, divergence constraints, historical or expert bounds, or another documented criterion. ARCCRAFT's proposed regime taxonomy can organize experiments, but the taxonomy contributes scientifically only when each regime is operationalized and its relationship to the baseline distribution is reported.

\subsection{Reverse stress testing and failure attribution}\label{sec:reverse-stress-testing-and-failure-attribution}

Forward stress testing begins with an adverse scenario and measures its consequences. Reverse stress testing begins with a failure condition and asks which changes to inputs or model distributions can produce that condition. Pesenti et al. \cite{Pesenti2019} provide a particularly relevant insurance application. Their reverse sensitivity framework specifies an output stress, derives a probability measure that satisfies the stress while minimizing divergence from the baseline, and evaluates how the distributions of input factors change. This design produces an interpretable ranking of factors associated with “breaking” a model without requiring repeated evaluation of an expensive aggregation function.

Kroell et al. \cite{Kroell2024} extend reverse stress concepts to dynamic loss models. They construct stressed probability measures for a compound Poisson process over time under terminal constraints, thereby showing that dynamic reverse stress testing is already an established research direction. ARCCRAFT should therefore not claim novelty for dynamic or reverse stress testing themselves. The relevant question is narrower: whether its procedural representation and Failure Atlas provide an auditable, useful link between simulated trajectories, first failure times, threshold conditions and replayable seeds that is not already supplied by the comparator methods.

Sensitivity analysis provides further context. Morris \cite{Morris1991} offers an efficient screening design for identifying influential inputs in computational experiments, whereas variance-based methods estimate main and interaction effects across the input space Saltelli et al. \cite{Saltelli2008,Saltelli2010}. These methods measure influence under specified distributions; they do not by themselves establish causality. A Failure Atlas that associates a failed seed with prominent variables should consequently use terms such as \emph{attribution}, \emph{association} or \emph{influence} unless it is tested against implanted causal mechanisms or another identification design.

This requirement suggests a falsifiable experiment. Synthetic worlds can be generated with registered faults whose responsible transition, dependency or event mechanism is known. ARCCRAFT and comparator diagnostics can then be evaluated blindly on their ability to detect the failure and recover its responsible mechanism. Precision, recall, false-discovery rate, detection delay and attribution rank would provide stronger evidence than selected examples of apparently informative failure traces.

\subsection{Synthetic insurance environments and open benchmarks}\label{sec:synthetic-insurance-environments-and-open-benchmarks}

Synthetic insurance data are not new. Gabrielli and Wuthrich \cite{Gabrielli2018} present an individual claims-history simulation machine calibrated on extensive real claim histories. Avanzi et al. \cite{Avanzi2021} introduce SynthETIC, an open simulator that gives users control over claim occurrence, notification, payment timing, inflation, dependencies and other realistic features. Avanzi et al. \cite{Avanzi2023} extend this ecosystem through SPLICE, which simulates paid and incurred claim histories and revisions of case estimates. These systems demonstrate that control over an insurance data-generating mechanism is valuable for benchmarking methods when confidential individual-claim data are unavailable.

Open artificial insurers provide an even closer architectural precedent. Wolf et al. \cite{Wolf2026} introduce openIRM as a publicly accessible internal risk model of an artificial life insurer for benchmarking nested-simulation and other actuarial methods under Solvency II. It integrates an economic scenario generator, asset-liability dynamics and capital estimation, and is calibrated over multiple market dates. This work shows that openness, time-dependent risk factors and reproducible benchmarking are already active contributions in actuarial science.

ARCCRAFT must consequently be distinguished by a specific combination of functions rather than by synthetic simulation alone. Its provisional distinction is the integration of versioned world specifications, explicit state transitions and events, separate stochastic streams, structural accept-repair-reject decisions, deterministic seed-level replay and a persistent failure-attribution record. Whether this integration yields scientific value remains an empirical question. It could improve diagnosis and governance without improving predictive accuracy; it could also increase computational cost or induce selection bias through validator decisions. The evaluation must allow each of these outcomes.

Synthetic evidence also has a fundamental limitation. A simulator can confirm that an analysis procedure recovers patterns deliberately embedded in the simulator, but it cannot independently demonstrate that those patterns describe the real insurance system. External portfolio evidence is therefore necessary. Synthetic fault injection can test failure detection and attribution, while independent observed outcomes must test predictive calibration and practical relevance. Neither form of evidence substitutes for the other.

\subsection{Reproducibility is not validation}\label{sec:reproducibility-is-not-validation}

Computational reproducibility requires enough information and artifacts to regenerate reported results. Peng \cite{Peng2011} and Stodden \cite{Stodden2015} treat accessible code, data and workflow information as central to reproducible computational research. Stodden et al. \cite{Stodden2016} further emphasize institutional practices that make computational methods inspectable and repeatable. In stochastic simulation, controlled random streams and substreams support designed comparisons and component isolation \cite{LEcuyer2002}. PCG64 provides an implementation choice for pseudo-random generation, but reporting a generator name and seed is insufficient unless software versions, stream construction and replay tolerances are also recorded.

Reproducibility nevertheless answers only whether an experiment can be repeated, not whether its model is empirically adequate. Sargent \cite{Sargent2013} distinguishes conceptual-model validity, model verification, operational validity and data validity. Kleijnen and Sargent \cite{Kleijnen2000} likewise distinguish fitting from validation and tie the required accuracy to the model's intended purpose. In actuarial practice, ASOP No. 56 requires attention to intended use, input reconciliation, formula and logic checks, sensitivity testing, output validation, limitations and model risk \cite{ASB2020}.

ARCCRAFT should therefore report at least three separate claims. \textbf{Software verification} concerns whether the coded engine implements its stated rules and reproduces registered runs. \textbf{Empirical validation} concerns whether output distributions are compatible with independent observations for their intended use. \textbf{Scenario plausibility} concerns whether adverse worlds remain defensible under empirical, regulatory or expert constraints. A successful replay establishes the first claim only. It cannot compensate for failed predictive coverage or an unjustified severe-plausible label.

The same separation applies to uncertainty. Der Kiureghian and Ditlevsen \cite{DerKiureghian2009} distinguish aleatory variability from epistemic uncertainty, while Kennedy and O'Hagan \cite{Kennedy2001} formalize the discrepancy between a computer model and the process it approximates. Barrieu and Scandolo \cite{Barrieu2015} show how model risk can be quantified relative to alternative models. These perspectives imply that one stochastic stream should not silently represent process randomness, parameter uncertainty, behavioural uncertainty and model-form uncertainty as if they were interchangeable. ARCCRAFT's independent streams may improve experimental control, but the semantic meaning and calibration of each source of randomness must be explicit.

\subsection{Predictive and tail evaluation}\label{sec:predictive-and-tail-evaluation}

Point accuracy alone is inadequate for a stochastic stress framework. Gneiting and Raftery \cite{Gneiting2007Scoring} show that proper scoring rules reward honest probabilistic forecasts, while Gneiting et al. \cite{Gneiting2007Calibration} distinguish calibration from sharpness. Czado et al. \cite{Czado2009} adapt predictive assessment to count data, which is directly relevant to insurance claim counts. A model can achieve a small error in its predictive mean and still be poorly calibrated if the observed outcome falls systematically outside its prediction intervals.

Tail evaluation requires similar care. Value-at-Risk exceedances can be assessed through unconditional coverage \cite{Kupiec1995} and conditional coverage \cite{Christoffersen1998}, while joint evaluation of Value-at-Risk and Expected Shortfall requires methods consistent with their statistical properties \cite{Nolde2017}. These tests require repeated observations. A single annual outcome above a 97.5\% predictive quantile is important negative evidence, but it is not enough to estimate a stable exceedance rate or to perform a well-powered backtest.

The present motor experiment illustrates why these distinctions matter. A retrospective projection calibrated on 2022-2023 produces a 2024 mean claim-count error of approximately -2.38\% and an incurred-cost error of approximately -4.86\%. Those central estimates appear close to the observed aggregates. However, the observed claim count, incurred cost and loss ratio all exceed their simulated 97.5\% quantiles. The experiment therefore cannot support complete predictive validation. It indicates that the central level is reasonably close while the reported predictive distribution understates the realized upper-tail outcome. Because 2024 has already been inspected, further tuning against that year would convert it from evaluation evidence into development data.

Distribution shift is one possible explanation, but it must be tested rather than presumed. Cauchois et al. \cite{Cauchois2024} show that predictive validity can deteriorate when training and evaluation distributions differ and develop prediction sets robust to bounded shifts. In the current application, changes in exposure mix, inflation, reporting practice, dependence or unmodelled portfolio characteristics could each contribute. The empirical design should use rolling-origin evaluation, an independent portfolio or a later untouched period, and it should report both central and distributional performance.

\subsection{Research gap and evaluation implications}\label{sec:research-gap-and-evaluation-implications}

The literature does not support describing ARCCRAFT as a new theory of actuarial stress testing. It supports a more precise gap. Established methods address dependent frequency-severity modelling, severe-plausible scenario selection, reverse stress testing, dynamic loss processes, synthetic claims, open artificial insurers, reproducible random streams, simulation validation and tail evaluation. What remains to be tested is whether integrating versioned procedural worlds, structural validation, controlled replay and seed-level failure attribution in one non-life workflow produces measurable benefits beyond those components and their conventional alternatives.

This gap leads directly to the research design. ARCCRAFT must be compared with deterministic projection, conventional frequency-severity Monte Carlo, a dependence-aware alternative and static stress scenarios under matched data and computational budgets. Its architectural components must be ablated. Validator repairs must be treated as distribution-changing interventions and evaluated accordingly. Failure detection and attribution must be tested against registered ground truth. Predictive performance must be assessed with proper scores, interval coverage and tail diagnostics, while reproducibility must be tested through exact or tolerance-based replay. Runtime, memory, rejection rates and effective accepted-sample size must be reported alongside any gain.

Under this design, ARCCRAFT can succeed in more than one defensible way. It may improve predictive distributions, detect failure configurations more efficiently, increase diagnostic accuracy, or strengthen reproducible model governance. It may also prove valuable mainly as an audit architecture without outperforming simpler predictive models. Conversely, if repairs distort the stress distribution, tail calibration remains poor, or the Failure Atlas does not recover known mechanisms, those findings would delimit the framework's value. A credible contribution depends on reporting these possibilities rather than defining success after observing the results.

\subsection{Testable expectations}\label{sec:testable-expectations}

The evaluation addresses the following expectations:

\begin{itemize}[leftmargin=*]
\item \textbf{H1:} a registered seed reproduces the same worlds, metrics, and failure labels in the inspected environment;
\item \textbf{H2:} named substreams isolate claims and behaviour from irrelevant consumption of the world stream;
\item \textbf{H3:} the validator changes the accepted scenario population and reported failures;
\item \textbf{H4:} multi-period state transitions change failure metrics relative to fixed validated states;
\item \textbf{H5:} policy-level models improve held-out predictive performance relative to pooled means;
\item \textbf{H6:} the Failure Atlas exposes replayable coded failure conditions, but causal diagnostic value remains unsupported without external or implanted ground truth.
\end{itemize}

H5 is evaluated dimension by dimension rather than through a post-hoc composite score.
